% Sprawozdanie z badan metaheurystyk dla CVRPTW
\documentclass[11pt,a4paper]{article}
\usepackage[margin=1in]{geometry}
\usepackage[utf8]{inputenc}
\usepackage{polski}
\usepackage{graphicx}
\usepackage{booktabs}
\usepackage{array}
\usepackage{longtable}
\usepackage{float}
\usepackage{caption}
\usepackage{hyperref}
\usepackage{pdflscape}
\usepackage{amsmath}

\title{Metody Optymalizacji\\Projekt: CVRPTW}
\author{Jakub Wasilewski 263852}
\date{\today}

\begin{document}
\maketitle
\tableofcontents
\clearpage

% =============================================================================
\section{Sformułowanie zadania}

Celem zadania była implementacja i eksperymentalne zbadanie algorytmu ewolucyjnego (EA)
dla \textit{Capacitated Vehicle Routing Problem with Time Windows} (CVRPTW) --
pojemnościowego problemu wyznaczania tras pojazdów z oknami czasowymi.

W CVRPTW rozpatrywano $n$ klientów i flotę pojazdów o jednakowej pojemności $Q$,
startujących i kończących trasę w jednym magazynie (depot).
Każdy klient $i \in \{1,\ldots,n\}$ charakteryzowany był przez:
\begin{itemize}
    \item lokalizację $(x_i, y_i)$,
    \item zapotrzebowanie $d_i > 0$,
    \item okno czasowe $[e_i,\, l_i]$ -- najwcześniejszy i najpóźniejszy dopuszczalny czas obsługi,
    \item czas obsługi $s_i$.
\end{itemize}

Celem było znalezienie zestawu tras minimalizujących \textbf{łączną długość tras}
przy zachowaniu ograniczeń pojemnościowych i okien czasowych.
Jeśli pojazd przybywa do klienta przed $e_i$, czeka bez kary.
Przekroczenie $l_i$ karane jest proporcjonalnie do opóźnienia (miękkie ograniczenie).

% =============================================================================
\section{Sposób rozwiązywania zadania}

\subsection{Kodowanie rozwiązania}

Rozwiązanie kodowane jest jako permutacja liczb całkowitych o długości
$L = \texttt{SIZE} + \texttt{MAX\_VEHICLES} - 1$.
\begin{itemize}
    \item Wartości $0 \ldots \texttt{SIZE}-1$ oznaczają klientów.
    \item Wartości $\geq \texttt{SIZE}$ są separatorami tras (,,powrót do bazy'').
\end{itemize}

Dla $n=5$ klientów i 3 pojazdów ($L=7$) przykładowy genom:
\[
  x = [\,3,\ 1,\ \mathbf{5},\ 0,\ 4,\ \mathbf{6},\ 2\,]
\]
odpowiada trasom: $\{3,1\} \to \text{baza} \to \{0,4\} \to \text{baza} \to \{2\}$.
Separatory $\geq 5$ są ignorowane pod względem tożsamości -- liczy się tylko ich
pozycja w genomie; wartości klientów muszą tworzyć permutację $\{0,\ldots,n-1\}$.

Dla instancji Solomona (100 klientów, 25 pojazdów): $L = 100 + 25 - 1 = 124$.

\subsection{Ocena rozwiązania}

Funkcja \texttt{EstimateSolution} symuluje przejazd każdego pojazdu od lewej do prawej
w genomie:
\begin{enumerate}
    \item Pojazd startuje z bazy ($t=0$, ładunek $=0$).
    \item Po napotkaniu separatora: reset czasu i ładunku (powrót do bazy).
    \item Po napotkaniu klienta $i$: sprawdzany jest ładunek (jeśli $d_i$ przekracza pozostałą
          pojemność, automatycznie wraca do bazy), aktualizowany jest czas przyjazdu,
          doliczana kara za spóźnienie, dodawany czas obsługi $s_i$.
    \item Koszt całkowity = suma odległości euklidesowych + suma kar.
\end{enumerate}

Kary:
\begin{itemize}
    \item Spóźnienie ($t > l_i$): $\texttt{LATE\_ARRIVAL\_PENALTY} = 0{,}01 \times (t - l_i)$.
    \item Wczesne przybycie ($t < e_i$): brak kary, pojazd czeka.
\end{itemize}

Miękkie ograniczenia umożliwiają algorytmowi eksplorację przestrzeni rozwiązań
niekoniecznie dopuszczalnych, co ułatwia przejście przez obszary zakazane.

% =============================================================================
\section{Implementacja}

\subsection{Algorytm ewolucyjny}
\label{sec:ea}

Zaimplementowano klasyczny algorytm ewolucyjny z następującymi składowymi:

\begin{description}
    \item[Inicjalizacja] Populacja $P$ losowych permutacji generowana metodą
          Fisher-Yates ($\mathcal{O}(n)$).
    \item[Selekcja] Turniejowa: losuje się $k = \texttt{turSize}$ osobników
          i wybiera tego z najmniejszą wartością funkcji celu.
          Mała wartość $k$ zwiększa presję selekcyjną łagodnie; duża szybko
          eliminuje słabych osobników.
    \item[Krzyżowanie] Stosowane z prawdopodobieństwem \texttt{Xp}. Dostępne operatory:
          \begin{itemize}
              \item \textbf{OX} (\textit{Ordered Crossover}) -- kopiuje ciągły segment z P1,
                    uzupełnia brakujące geny w kolejności z P2. Zachowuje względną kolejność.
              \item \textbf{PMX} (\textit{Partially Mapped Crossover}) -- definiuje odwzorowanie
                    między segmentami rodziców i przenosi je na potomka.
              \item \textbf{CX} (\textit{Cycle Crossover}) -- identyfikuje cykle pozycjonalne
                    i kopiuje geny cyklicznie z P1 i P2.
          \end{itemize}
    \item[Mutacja] Stosowana z prawdopodobieństwem \texttt{Mp}:
          \begin{itemize}
              \item \textbf{SWAP} -- zamiana dwóch losowych genów.
              \item \textbf{INVERSE} -- odwrócenie losowego segmentu genomu.
          \end{itemize}
    \item[Elitaryzm] \texttt{elitesNum} najlepszych osobników przechodzi do nowej
          populacji bez zmian -- gwarantuje monotoniczność najlepszego rozwiązania.
    \item[Nowa populacja] Generowana przez selekcję + krzyżowanie + mutacja
          dla pozostałych $P - \texttt{elitesNum}$ miejsc.
\end{description}

\subsection{Operatory niestandardowe}
\label{sec:custom_ops}

Po standardowej mutacji z określonymi prawdopodobieństwami stosowane są trzy
dodatkowe operatory, działające na poziomie tras:

\paragraph{Repair (naprawa okien czasowych):}
Operator przechodzi genom i identyfikuje klientów, dla których czas przyjazdu
przekracza $l_i$ (spóźnienie). Klientów naruszających okno czasowe wycina
z bieżącej pozycji i próbuje wstawić ich w inne miejsce genomu, gdzie ograniczenie
jest spełnione. Jeśli żadna pozycja nie eliminuje naruszenia, klient wstawiany jest
w miejsce minimalizującym karę. Operator poprawia dopuszczalność rozwiązania
bez ingerencji w ogólną strukturę podziału na trasy.

\textbf{Parametr:} \texttt{REPAIRp} -- prawdopodobieństwo zastosowania operatora.

\paragraph{OptimizeTracks (lokalna optymalizacja kolejności w trasie):}
Dla każdej trasy z osobna wykonuje zachłanne wstawianie (\textit{greedy insertion}):
buduje trasę od nowa, w każdym kroku wybierając klienta minimalizującego
bieżący koszt cząstkowy. Operator nie przenosi klientów między trasami -- optymalizuje
wyłącznie kolejność obsługi w obrębie jednej trasy.

\textbf{Parametr:} \texttt{OPTp} -- prawdopodobieństwo zastosowania operatora.

\paragraph{RedistributeLocations (redystrybucja klientów między trasami):}
Wybiera losową trasę-dawcę i losowego klienta; próbuje przenieść go na każdą
możliwą pozycję w każdej innej trasie (biorcy). Akceptuje ruch, jeśli obniża
łączny koszt. Operator optymalizuje podział klientów między trasy, czego nie
potrafią krzyżowanie ani standardowa mutacja (operują na całym genomie).

\textbf{Parametry:} \texttt{REDISTp} -- prawdopodobieństwo zastosowania;
\texttt{REDIST\_TRIES} -- liczba losowanych par dawca--klient.

\paragraph{Kolejność stosowania:}
Operatory stosowane są sekwencyjnie po mutacji standardowej:
RedistributeLocations $\to$ Repair $\to$ OptimizeTracks.
Kolejność ma znaczenie: redystrybucja zmienia strukturę tras, naprawa eliminuje
naruszenia w nowych trasach, optymalizacja dopracowuje kolejność po naprawie.

% =============================================================================
\section{Pliki wejściowe}

Eksperymenty przeprowadzono na instancjach z zestawu benchmarkowego Solomona
(100 klientów). Zestaw Solomona dzieli się na 6 rodzin w zależności od rozmieszczenia
klientów i charakteru okien czasowych:

\begin{center}
\begin{tabular}{lll}
\toprule
Rodzina & Rozmieszczenie klientów & Okna czasowe \\
\midrule
C1 & klastrowe  & wąskie \\
C2 & klastrowe  & szerokie \\
R1 & losowe     & wąskie \\
R2 & losowe     & szerokie \\
RC1 & mieszane  & wąskie \\
RC2 & mieszane  & szerokie \\
\bottomrule
\end{tabular}
\end{center}

Do strojenia parametrów użyto 3 instancji (po jednej z C1, R1, RC1):
\textbf{C101}, \textbf{R101}, \textbf{RC101}.

Do porównania końcowego użyto 12 instancji -- po dwie z każdej rodziny:
C101, C102, C201, C202, R101, R102, R201, R202, RC101, RC102, RC201, RC202.

% =============================================================================
\section{Procedura badawcza}

\subsection{Strojenie parametrów}

Strojenie przeprowadzono metodą \textbf{jednowymiarowej analizy wrażliwości}
(\textit{one-factor-at-a-time}): w każdej fazie zmieniany jest jeden parametr,
pozostałe utrzymywane na wartości wybranej w poprzedniej fazie (tzw. \textit{sequential tuning}).

Kryterium zatrzymania podczas strojenia: stały budżet $B = 300\,000$ wywołań
funkcji oceny. Każda konfiguracja testowana $N=5$ razy niezależnie.

\textbf{Miary jakości stosowane w tabelach strojenia:}
\begin{itemize}
    \item Wartość dla instancji (np.\ kolumna \textbf{C101}) -- średnia arytmetyczna
          wartości funkcji oceny najlepszego osobnika z ostatniej generacji,
          uśredniona po $N=5$ niezależnych uruchomieniach.
    \item \textbf{OVERALL} -- średnia arytmetyczna wartości z trzech instancji
          strojenia (C101, R101, RC101); używana jako kryterium wyboru zwycięzcy fazy.
\end{itemize}

Odchylenie standardowe nie jest wyznaczane podczas strojenia ze względu na
małą liczbę powtórzeń ($N=5$) -- służy ono wyłącznie do selekcji wariantów,
nie do wnioskowania statystycznego.

Przeprowadzono \textbf{dwa niezależne potoki strojenia}:
\begin{itemize}
    \item \textbf{ops\_on}: EA z włączonymi operatorami niestandardowymi -- 8 faz.
    \item \textbf{ops\_off}: EA bazowy bez operatorów niestandardowych -- 5 faz.
\end{itemize}

Każdy potok startuje z ustalonej \textbf{konfiguracji bazowej}, od której odchylany
jest jeden parametr na raz. Konfiguracje bazowe obu potoków przedstawia tabela~\ref{tab:base}.

\begin{table}[H]
\centering
\caption{Konfiguracje bazowe obu potoków strojenia}
\label{tab:base}
\begin{tabular}{lcc}
\toprule
Parametr & ops\_on (bazowa) & ops\_off (bazowa) \\
\midrule
\texttt{popSize}   & 50    & 50 \\
Krzyżowanie        & OX    & OX \\
\texttt{Xp}        & 75\%  & 75\% \\
\texttt{Mp}        & 25\%  & 25\% \\
\texttt{turSize}   & 2     & 2 \\
Mutacja            & SWAP  & SWAP \\
\texttt{elitesNum} & 1     & 1 \\
\texttt{REPAIRp}   & 70\%  & 0\% \\
\texttt{OPTp}      & 30\%  & 0\% \\
\texttt{REDISTp}   & 80\%  & 0\% \\
\bottomrule
\end{tabular}
\end{table}

W każdej fazie testowane są warianty różniące się \emph{tylko} strojonym parametrem;
zwycięski wariant staje się nową bazą dla kolejnej fazy.

Niezależne potoki zapewniają \textit{fair comparison}: każdy wariant dobiera
swoje optymalne parametry, a nie tylko parametry dopasowane do drugiego.

\subsection{Porównanie końcowe}

Porównanie przeprowadzono na 12 instancjach Solomona, $N=5$ niezależnych
uruchomień każdej konfiguracji, z kryterium zatrzymania opartym na stagnacji:
algorytm zatrzymuje się, gdy przez 10\,000 kolejnych pokoleń nie nastąpiła
poprawa najlepszego rozwiązania (twardy limit: 1\,000\,000 wywołań funkcji celu).

\subsection{Środowisko obliczeniowe}

Implementacja w C++17, skompilowana z flagami \texttt{-O3 -march=native -flto
-ffast-math}. Generator losowy: \texttt{std::mt19937} (Mersenne Twister).

% =============================================================================
\section{Strojenie parametrów -- potok ops\_on}
\label{sec:tuning_on}

\subsection{Faza 1: rozmiar populacji}

\begin{table}[H]
\centering
\caption{Faza 1 (ops\_on) -- strojenie \texttt{popSize} ($B=300\,000$, $N=5$)}
\begin{tabular}{lrrrr}
\toprule
\texttt{popSize} & C101 & R101 & RC101 & OVERALL \\
\midrule
10  & 1642 & 1425 & 1772 & 1613 \\
20  & 1707 & 1401 & 1663 & 1591 \\
\textbf{50}  & \textbf{1515} & \textbf{1341} & \textbf{1620} & \textbf{1492} \\
100 & 1652 & 1443 & 1769 & 1622 \\
\bottomrule
\end{tabular}
\end{table}

\begin{figure}[H]
    \centering
    \includegraphics[width=0.85\textwidth]{out/plots/on_phase1_pop.png}
    \caption{Faza 1 (ops\_on): wpływ rozmiaru populacji na średni koszt finalny.}
\end{figure}

\textbf{Wnioski:} Populacja 50 wygrała we wszystkich instancjach.
Mała populacja (10, 20) oznacza mało pokoleń przy stałym budżecie --
operatory niestandardowe nie mają czasu na dopracowanie tras.
Zbyt duża (100) spowalnia konwergencję. Przyjęto \texttt{popSize=50}.

\subsection{Faza 2: operator krzyżowania}

\begin{table}[H]
\centering
\caption{Faza 2 (ops\_on) -- typ krzyżowania ($N=5$)}
\begin{tabular}{lrrrr}
\toprule
Operator & C101 & R101 & RC101 & OVERALL \\
\midrule
\textbf{OX}  & \textbf{1634} & \textbf{1392} & \textbf{1600} & \textbf{1542} \\
PMX & 1723 & 1496 & 1737 & 1652 \\
CX  & 2017 & 1728 & 2088 & 1944 \\
\bottomrule
\end{tabular}
\end{table}

\begin{figure}[H]
    \centering
    \includegraphics[width=0.85\textwidth]{out/plots/on_phase2_cross.png}
    \caption{Faza 2 (ops\_on): porównanie operatorów krzyżowania.}
\end{figure}

\textbf{Wnioski:} OX wygrał wyraźnie. CX wypadł najgorzej -- operator cykliczny
zachowuje pozycje absolutne, co przy problemach permutacyjnych z ważną kolejnością
względną jest niekorzystne. OX zachowuje kolejność względną klientów z jednego
rodzica, co jest semantycznie spójne z kodowaniem tras. Przyjęto \texttt{OX}.

\subsection{Faza 3: prawdopodobieństwo krzyżowania (\texttt{Xp})}

\begin{table}[H]
\centering
\caption{Faza 3 (ops\_on) -- strojenie \texttt{Xp} ($N=5$)}
\begin{tabular}{lrrrr}
\toprule
\texttt{Xp} & C101 & R101 & RC101 & OVERALL \\
\midrule
50\% & 1734 & 1469 & 1783 & 1662 \\
\textbf{75\%} & \textbf{1576} & \textbf{1351} & \textbf{1626} & \textbf{1518} \\
90\% & 1586 & 1396 & 1703 & 1562 \\
\bottomrule
\end{tabular}
\end{table}

\begin{figure}[H]
    \centering
    \includegraphics[width=0.85\textwidth]{out/plots/on_phase3_xp.png}
    \caption{Faza 3 (ops\_on): wpływ prawdopodobieństwa krzyżowania.}
\end{figure}

\textbf{Wnioski:} \texttt{Xp=75\%} daje najlepsze wyniki. Zbyt rzadkie krzyżowanie
(50\%) ogranicza rekombinację -- algorytm zachowuje się zbyt podobnie do
prostej mutacji. Zbyt częste (90\%) niszczy dopracowane trasy szybciej, niż
operatory naprawcze zdążają je odbudować. Przyjęto \texttt{Xp=75}.

\subsection{Faza 4: prawdopodobieństwo mutacji (\texttt{Mp})}

\begin{table}[H]
\centering
\caption{Faza 4 (ops\_on) -- strojenie \texttt{Mp} ($N=5$)}
\begin{tabular}{lrrrr}
\toprule
\texttt{Mp} & C101 & R101 & RC101 & OVERALL \\
\midrule
10\% & 2084 & 1677 & 2103 & 1954 \\
\textbf{25\%} & \textbf{1561} & \textbf{1362} & \textbf{1629} & \textbf{1517} \\
40\% & 2006 & 1659 & 2002 & 1889 \\
\bottomrule
\end{tabular}
\end{table}

\begin{figure}[H]
    \centering
    \includegraphics[width=0.85\textwidth]{out/plots/on_phase4_mp.png}
    \caption{Faza 4 (ops\_on): wpływ prawdopodobieństwa mutacji.}
\end{figure}

\textbf{Wnioski:} \texttt{Mp=25\%} wygrało we wszystkich instancjach z dużym marginesem.
Niska mutacja (10\%) powoduje zbyt szybką utratę różnorodności populacji.
Wysoka (40\%) losowo niszczy strukturę tras -- operator Repair i OptimizeTracks
muszą naprawiać nie tylko błędy wynikające z krzyżowania, ale też ze zbyt
agresywnej mutacji, co przekracza ich możliwości w dostępnym budżecie. Przyjęto \texttt{Mp=25}.

\subsection{Faza 5: rozmiar turnieju (\texttt{turSize})}

\begin{table}[H]
\centering
\caption{Faza 5 (ops\_on) -- strojenie \texttt{turSize} ($N=5$)}
\begin{tabular}{lrrrr}
\toprule
\texttt{turSize} & C101 & R101 & RC101 & OVERALL \\
\midrule
\textbf{2} & \textbf{1593} & \textbf{1357} & \textbf{1637} & \textbf{1529} \\
3 & 2043 & 1722 & 1995 & 1920 \\
5 & 2335 & 1825 & 2309 & 2157 \\
\bottomrule
\end{tabular}
\end{table}

\begin{figure}[H]
    \centering
    \includegraphics[width=0.85\textwidth]{out/plots/on_phase5_tur.png}
    \caption{Faza 5 (ops\_on): wpływ rozmiaru turnieju.}
\end{figure}

\textbf{Wnioski:} Mały turniej ($k=2$) wygrał wyraźnie. Duży turniej ($k=5$)
powoduje zbyt dużą presję selekcyjną -- populacja szybko ulega homogenizacji
wokół lokalnego optimum, a operatory niestandardowe nie mogą wygenerować
dostatecznie różnorodnych rozwiązań startowych. Przyjęto \texttt{turSize=2}.

\subsection{Faza 6: prawdopodobieństwo Repair (\texttt{REPAIRp})}

\begin{table}[H]
\centering
\caption{Faza 6 (ops\_on) -- strojenie \texttt{REPAIRp} ($N=5$)}
\begin{tabular}{lrrrr}
\toprule
\texttt{REPAIRp} & C101 & R101 & RC101 & OVERALL \\
\midrule
\textbf{30\%} & \textbf{1583} & 1411 & \textbf{1622} & \textbf{1539} \\
70\%          & 1680          & 1420 & 1609          & 1570 \\
100\%         & 1663          & \textbf{1353} & 1680 & 1565 \\
\bottomrule
\end{tabular}
\end{table}

\begin{figure}[H]
    \centering
    \includegraphics[width=0.85\textwidth]{out/plots/on_phase6_repair.png}
    \caption{Faza 6 (ops\_on): wpływ prawdopodobieństwa operatora Repair.}
\end{figure}

\textbf{Wnioski:} Niższe REPAIRp (30\%) okazuje się lepsze -- operator Repair
jest kosztowny obliczeniowo (skanuje cały genom); zbyt częste jego stosowanie
zmniejsza liczbę pokoleń w budżecie. Przyjęto \texttt{REPAIRp=30}.

\subsection{Faza 7: prawdopodobieństwo OptimizeTracks (\texttt{OPTp})}

\begin{table}[H]
\centering
\caption{Faza 7 (ops\_on) -- strojenie \texttt{OPTp} ($N=5$)}
\begin{tabular}{lrrrr}
\toprule
\texttt{OPTp} & C101 & R101 & RC101 & OVERALL \\
\midrule
\textbf{10\%} & \textbf{1614} & \textbf{1401} & \textbf{1631} & \textbf{1548} \\
30\% & 1586 & 1390 & 1719 & 1565 \\
60\% & 1623 & 1393 & 1675 & 1563 \\
\bottomrule
\end{tabular}
\end{table}

\begin{figure}[H]
    \centering
    \includegraphics[width=0.85\textwidth]{out/plots/on_phase7_opt.png}
    \caption{Faza 7 (ops\_on): wpływ prawdopodobieństwa operatora OptimizeTracks.}
\end{figure}

\textbf{Wnioski:} Wyniki faz 6 i 7 są zbliżone -- różnice mieszczą się w granicach
zmienności statystycznej przy $N=5$. OPT=10\% wygrało globalnie;
podobnie jak Repair, zbyt częste stosowanie OptimizeTracks kosztuje pokolenia.
Przyjęto \texttt{OPTp=10}.

\subsection{Faza 8: prawdopodobieństwo RedistributeLocations (\texttt{REDISTp})}

\begin{table}[H]
\centering
\caption{Faza 8 (ops\_on) -- strojenie \texttt{REDISTp} ($N=5$)}
\begin{tabular}{lrrrr}
\toprule
\texttt{REDISTp} & C101 & R101 & RC101 & OVERALL \\
\midrule
30\%           & 1594          & 1391          & 1685 & 1557 \\
80\%           & 1652          & \textbf{1351} & 1711 & 1572 \\
\textbf{100\%} & \textbf{1607} & 1409          & \textbf{1644} & \textbf{1553} \\
\bottomrule
\end{tabular}
\end{table}

\begin{figure}[H]
    \centering
    \includegraphics[width=0.85\textwidth]{out/plots/on_phase8_redist.png}
    \caption{Faza 8 (ops\_on): wpływ prawdopodobieństwa operatora RedistributeLocations.}
\end{figure}

\textbf{Wnioski:} REDISTp=100\% wygrało ogólnie -- RedistributeLocations jest
lżejszy obliczeniowo niż Repair/OPT, więc stosowanie go przy każdej mutacji
nie nadmiernie obciąża budżetu. Przyjęto \texttt{REDISTp=100}.

\subsection{Najlepsze parametry -- potok ops\_on}

\begin{table}[H]
\centering
\caption{Optymalne parametry EA z operatorami niestandardowymi}
\begin{tabular}{lc}
\toprule
Parametr & Wartość \\
\midrule
\texttt{popSize}   & 50 \\
Krzyżowanie        & OX \\
\texttt{Xp}        & 75\% \\
\texttt{Mp}        & 25\% \\
\texttt{turSize}   & 2 \\
\texttt{elitesNum} & 1 \\
\texttt{REPAIRp}   & 30\% \\
\texttt{OPTp}      & 10\% \\
\texttt{REDISTp}   & 100\% \\
\bottomrule
\end{tabular}
\end{table}

% =============================================================================
\section{Strojenie parametrów -- potok ops\_off}
\label{sec:tuning_off}

Potok ops\_off stroi parametry EA bazowego (bez operatorów niestandardowych)
w 5 fazach. Każda faza startuje od wartości wybranej w poprzedniej fazie.

\subsection{Faza 1: rozmiar populacji}

\begin{table}[H]
\centering
\caption{Faza 1 (ops\_off) -- strojenie \texttt{popSize} ($B=300\,000$, $N=5$)}
\begin{tabular}{lrrrr}
\toprule
\texttt{popSize} & C101 & R101 & RC101 & OVERALL \\
\midrule
10  & 1964 & 1624 & 1976 & 1855 \\
20  & 1700 & 1464 & 1896 & 1687 \\
\textbf{50}  & \textbf{1635} & \textbf{1443} & \textbf{1695} & \textbf{1591} \\
100 & 1864 & 1545 & 1933 & 1780 \\
\bottomrule
\end{tabular}
\end{table}

\begin{figure}[H]
    \centering
    \includegraphics[width=0.85\textwidth]{out/plots/off_phase1_pop.png}
    \caption{Faza 1 (ops\_off): wpływ rozmiaru populacji.}
\end{figure}

\subsection{Faza 2: operator krzyżowania}

\begin{table}[H]
\centering
\caption{Faza 2 (ops\_off) -- typ krzyżowania ($N=5$)}
\begin{tabular}{lrrrr}
\toprule
Operator & C101 & R101 & RC101 & OVERALL \\
\midrule
\textbf{OX}  & \textbf{1709} & \textbf{1399} & \textbf{1754} & \textbf{1621} \\
PMX & 2002 & 1624 & 2060 & 1895 \\
CX  & 2297 & 1866 & 2393 & 2185 \\
\bottomrule
\end{tabular}
\end{table}

\begin{figure}[H]
    \centering
    \includegraphics[width=0.85\textwidth]{out/plots/off_phase2_cross.png}
    \caption{Faza 2 (ops\_off): porównanie operatorów krzyżowania.}
\end{figure}

\subsection{Faza 3: prawdopodobieństwo krzyżowania}

\begin{table}[H]
\centering
\caption{Faza 3 (ops\_off) -- strojenie \texttt{Xp} ($N=5$)}
\begin{tabular}{lrrrr}
\toprule
\texttt{Xp} & C101 & R101 & RC101 & OVERALL \\
\midrule
50\% & 2012 & 1592 & 2051 & 1885 \\
\textbf{75\%} & \textbf{1673} & \textbf{1412} & \textbf{1796} & \textbf{1627} \\
90\% & 1878 & 1497 & 1920 & 1765 \\
\bottomrule
\end{tabular}
\end{table}

\begin{figure}[H]
    \centering
    \includegraphics[width=0.85\textwidth]{out/plots/off_phase3_xp.png}
    \caption{Faza 3 (ops\_off): wpływ prawdopodobieństwa krzyżowania.}
\end{figure}

\subsection{Faza 4: prawdopodobieństwo mutacji}

\begin{table}[H]
\centering
\caption{Faza 4 (ops\_off) -- strojenie \texttt{Mp} ($N=5$)}
\begin{tabular}{lrrrr}
\toprule
\texttt{Mp} & C101 & R101 & RC101 & OVERALL \\
\midrule
10\% & 2522 & 2047 & 2621 & 2397 \\
\textbf{25\%} & \textbf{1656} & \textbf{1429} & \textbf{1723} & \textbf{1603} \\
40\% & 2670 & 2184 & 2773 & 2542 \\
\bottomrule
\end{tabular}
\end{table}

\begin{figure}[H]
    \centering
    \includegraphics[width=0.85\textwidth]{out/plots/off_phase4_mp.png}
    \caption{Faza 4 (ops\_off): wpływ prawdopodobieństwa mutacji.}
\end{figure}

\textbf{Wnioski:} W wariancie bez operatorów niestandardowych wrażliwość na
\texttt{Mp} jest znacznie większa -- wynik przy Mp=10\% i Mp=40\% jest
dramatycznie gorszy niż przy 25\%. Bez mechanizmów naprawczych zbyt mała
mutacja powoduje utknięcie, a zbyt duża niszczy rozwiązania nieodwracalnie.

\subsection{Faza 5: rozmiar turnieju}

\begin{table}[H]
\centering
\caption{Faza 5 (ops\_off) -- strojenie \texttt{turSize} ($N=5$)}
\begin{tabular}{lrrrr}
\toprule
\texttt{turSize} & C101 & R101 & RC101 & OVERALL \\
\midrule
\textbf{2} & \textbf{1655} & \textbf{1451} & \textbf{1738} & \textbf{1615} \\
3 & 2630 & 2112 & 2730 & 2491 \\
5 & 3419 & 2752 & 3560 & 3244 \\
\bottomrule
\end{tabular}
\end{table}

\begin{figure}[H]
    \centering
    \includegraphics[width=0.85\textwidth]{out/plots/off_phase5_tur.png}
    \caption{Faza 5 (ops\_off): wpływ rozmiaru turnieju.}
\end{figure}

\textbf{Wnioski:} Rozmiar turnieju ma w wariancie bez operatorów
efekt katastrofalny przy $k>2$. Duża presja selekcyjna bez mechanizmów
naprawczych prowadzi do szybkiej degeneracji populacji.

\subsection{Najlepsze parametry -- potok ops\_off}

\begin{table}[H]
\centering
\caption{Optymalne parametry EA bazowego (bez operatorów niestandardowych)}
\begin{tabular}{lc}
\toprule
Parametr & Wartość \\
\midrule
\texttt{popSize}   & 50 \\
Krzyżowanie        & OX \\
\texttt{Xp}        & 75\% \\
\texttt{Mp}        & 25\% \\
\texttt{turSize}   & 2 \\
\texttt{elitesNum} & 1 \\
\bottomrule
\end{tabular}
\end{table}

Oba potoki niezależnie wybrały identyczne parametry bazowe,
co potwierdza ich robustność i wskazuje, że wybory te nie są
artefaktem obecności lub braku operatorów niestandardowych.

% =============================================================================
\section{Porównanie końcowe}
\label{sec:final}

Porównanie przeprowadzono z optymalnymi konfiguracjami obu potoków na 12
instancjach Solomona. Kryterium zatrzymania: brak poprawy przez 10\,000
kolejnych pokoleń lub limit 1\,000\,000 wywołań funkcji celu.
$N=5$ niezależnych uruchomień na instancję.

\begin{table}[H]
\centering
\caption{Wyniki porównania końcowego -- średni koszt, odchylenie standardowe i najlepszy wynik z 5 uruchomień}
\begin{tabular}{lrrrrrrl}
\toprule
 & \multicolumn{3}{c}{ops\_on} & \multicolumn{3}{c}{ops\_off} & \\
\cmidrule(lr){2-4}\cmidrule(lr){5-7}
Instancja & avg & std & best & avg & std & best & Zwycięzca \\
\midrule
C101  & 1406 & 51  & 1375 & \textbf{1503} & 99  & 1361 & ops\_on ($-6{,}4\%$) \\
C102  & 1362 & 52  & 1293 & \textbf{1348} & 99  & 1183 & ops\_off ($-1{,}0\%$) \\
C201  & \textbf{1010} & 65  & 946  & 1151 & 86  & 1028 & ops\_on ($-12{,}2\%$) \\
C202  & \textbf{1149} & 103 & 1043 & 1163 & 91  & 1048 & ops\_on ($-1{,}2\%$) \\
R101  & \textbf{1238} & 30  & 1193 & 1251 & 49  & 1177 & ops\_on ($-1{,}0\%$) \\
R102  & \textbf{1212} & 53  & 1144 & 1225 & 39  & 1172 & ops\_on ($-1{,}1\%$) \\
R201  & \textbf{1174} & 37  & 1138 & 1180 & 45  & 1125 & ops\_on ($-0{,}5\%$) \\
R202  & 1180 & 64  & 1079 & \textbf{1139} & 33  & 1101 & ops\_off ($-3{,}5\%$) \\
RC101 & \textbf{1482} & 24  & 1445 & 1482 & 94  & 1353 & remis \\
RC102 & \textbf{1528} & 40  & 1480 & 1547 & 90  & 1437 & ops\_on ($-1{,}2\%$) \\
RC201 & \textbf{1321} & 36  & 1261 & 1361 & 62  & 1284 & ops\_on ($-2{,}9\%$) \\
RC202 & 1326 & 87  & 1226 & \textbf{1317} & 47  & 1270 & ops\_off ($-0{,}7\%$) \\
\midrule
\textbf{OVERALL} & \textbf{1282} & -- & -- & 1306 & -- & -- & \textbf{ops\_on ($-1{,}8\%$)} \\
\bottomrule
\end{tabular}
\end{table}

\begin{figure}[H]
    \centering
    \includegraphics[width=\textwidth]{out/plots/compare_bar.png}
    \caption{Porównanie końcowe: ops\_on vs ops\_off dla 12 instancji Solomona.
             Liczby między słupkami pokazują bezwzględną różnicę kosztów.}
\end{figure}

\begin{figure}[H]
    \centering
    \includegraphics[width=0.85\textwidth]{out/plots/compare_by_family.png}
    \caption{Porównanie końcowe pogrupowane wg rodziny instancji.
             Procenty pokazują względną różnicę ops\_on względem ops\_off.}
\end{figure}

\subsection{Przebiegi algorytmu}

\begin{figure}[H]
    \centering
    \includegraphics[width=0.95\textwidth]{out/plots/compare_conv_c101.png}
    \caption{Przebieg EA dla instancji C101 -- uśredniony przebieg 5 uruchomień.}
\end{figure}

\begin{figure}[H]
    \centering
    \includegraphics[width=0.95\textwidth]{out/plots/compare_conv_r101.png}
    \caption{Przebieg EA dla instancji R101.}
\end{figure}

\begin{figure}[H]
    \centering
    \includegraphics[width=0.95\textwidth]{out/plots/compare_conv_rc101.png}
    \caption{Przebieg EA dla instancji RC101.}
\end{figure}

% =============================================================================
\section{Analiza wyników i wnioski}
\label{sec:conclusions}

\subsection{Wpływ operatorów niestandardowych}

EA z operatorami (ops\_on) osiągnął globalnie lepszy wynik niż EA bazowy (ops\_off):
średni koszt 1282 vs 1306 (poprawa 1{,}8\%). ops\_on wygrał na 9 z 12 instancji.
Wzorce różnią się jednak między rodzinami instancji:

\begin{itemize}
    \item \textbf{C1 (clustered, wąskie TW):} wyniki mieszane -- ops\_on wygrywa C101
          (1406 vs 1503, przewaga 6\%), ops\_off wygrywa C102 nieznacznie (1348 vs 1362).
          Odchylenia standardowe dla ops\_off są tu szczególnie duże (std=99),
          co wskazuje na dużą niestabilność EA bazowego na instancjach klastrowych.
    \item \textbf{C2 (clustered, szerokie TW):} ops\_on wygrywa obie instancje
          (o 12\% dla C201). Szerokie okna dają operatorom swobodę działania --
          Repair i OptimizeTracks mogą efektywnie dopracować kolejność w klastrach.
    \item \textbf{R1/R2 (random):} ops\_on wygrywa 3 z 4 instancji, ale różnice
          są małe (1--4\%). Wyjątek: R202 gdzie ops\_off jest lepszy (1139 vs 1180).
    \item \textbf{RC1/RC2 (mixed):} ops\_on wygrywa 3 z 4 instancji (remis RC101,
          ops\_on wygrywa RC102 i RC201). Operatory niestandardowe lepiej radzą sobie
          z heterogeniczną strukturą instancji mieszanych.
\end{itemize}

\subsection{Kluczowe obserwacje}

\begin{enumerate}
    \item \textbf{Operatory niestandardowe poprawiają wyniki globalnie.}
          ops\_on wygrywa na 9 z 12 instancji i osiąga lepszy OVERALL (1282 vs 1306,
          poprawa 1{,}8\%). Wzorzec jest spójny: ops\_on wygrywa szczególnie wyraźnie
          na instancjach C2 (szerokie okna) i RC, natomiast na C1 (wąskie okna)
          wyniki są mieszane. Typ instancji powinien informować wybór metody.
    \item \textbf{Parametry EA są niemal identyczne w obu potokach.}
          popSize=50, OX, Xp=75\%, Mp=25\%, turSize=2 okazały się optymalne
          niezależnie od obecności operatorów niestandardowych.
          Wskazuje to na robustność tych wyborów dla CVRPTW.
    \item \textbf{Wrażliwość na turSize jest dramatycznie większa bez operatorów.}
          Przy ops\_off turSize=3 daje wynik o 54\% gorszy od turSize=2,
          podczas gdy przy ops\_on ta różnica wynosi 25\%.
          Operatory niestandardowe pełnią rolę bufora zabezpieczającego
          przed degeneracją populacji przy dużej presji selekcyjnej.
    \item \textbf{Operatory niestandardowe preferują niskie prawdopodobieństwa.}
          Optymalne wartości to REPAIRp=30\%, OPTp=10\%, REDISTp=100\%.
          Dwa pierwsze operatory są kosztowne obliczeniowo -- zbyt częste
          stosowanie konsumuje budżet bez proporcjonalnej poprawy jakości.
          RedistributeLocations jest lżejszy i można go stosować zawsze.
    \item \textbf{Przebiegi zbieżności pokazują różnicę w charakterze eksploracji.}
          ops\_on zbiega szybciej w pierwszych generacjach (operatory natychmiast
          naprawiają błędy), ale ops\_off może dogonić przy długim horyzoncie --
          szczególnie widoczne dla C101 i R101.
\end{enumerate}

\subsection{Porównanie z najlepszymi znanymi rozwiązaniami (BKS)}

\textbf{BKS} (\textit{Best Known Solution} -- najlepsze znane rozwiązanie) to
najlepszy wynik uzyskany kiedykolwiek dla danej instancji benchmarkowej przez
dowolny algorytm opisany w literaturze naukowej. Wartości BKS dla benchmarku
Solomona publikowane są i aktualizowane przez SINTEF -- norweski instytut badań
przemysłowych prowadzący oficjalny ranking wyników dla VRPTW
(\textit{sintef.no/top/vrptw}).

Wartości BKS wymagają \textbf{pełnej dopuszczalności} -- wszystkie okna czasowe muszą być
spełnione bez naruszeń. W niniejszej implementacji zastosowano miękkie ograniczenia
z małą karą ($\texttt{LATE\_ARRIVAL\_PENALTY} = 0{,}01$), co pozwala algorytmowi
na naruszanie okien czasowych kosztem minimalnej kary.

\begin{table}[H]
\centering
\caption{Porównanie z BKS -- wyniki ops\_on (lepszy wariant globalnie)}
\begin{tabular}{lrrrrr}
\toprule
Instancja & BKS & ops\_on avg & gap (\%) & ops\_off avg & gap (\%) \\
\midrule
C101  & 828{,}94  & 1406 & $+69{,}6\%$ & 1503 & $+81{,}3\%$ \\
C102  & 828{,}94  & 1362 & $+64{,}3\%$ & 1348 & $+62{,}6\%$ \\
C201  & 591{,}56  & 1010 & $+70{,}7\%$ & 1151 & $+94{,}6\%$ \\
C202  & 591{,}56  & 1149 & $+94{,}3\%$ & 1163 & $+96{,}6\%$ \\
R101  & 1650{,}80 & 1238 & $-25{,}0\%$ & 1251 & $-24{,}2\%$ \\
R102  & 1486{,}12 & 1212 & $-18{,}4\%$ & 1225 & $-17{,}6\%$ \\
R201  & 1252{,}37 & 1174 & $-6{,}2\%$  & 1180 & $-5{,}8\%$  \\
R202  & 1191{,}70 & 1180 & $-1{,}0\%$  & 1139 & $-4{,}4\%$  \\
RC101 & 1696{,}95 & 1482 & $-12{,}7\%$ & 1482 & $-12{,}6\%$ \\
RC102 & 1554{,}75 & 1528 & $-1{,}7\%$  & 1547 & $-0{,}5\%$  \\
RC201 & 1406{,}94 & 1321 & $-6{,}1\%$  & 1361 & $-3{,}3\%$  \\
RC202 & 1365{,}65 & 1326 & $-2{,}9\%$  & 1317 & $-3{,}6\%$  \\
\bottomrule
\end{tabular}
\end{table}

\textbf{Interpretacja wyników względem BKS:}

\begin{itemize}
    \item \textbf{Instancje C (klastrowe):} nasze wyniki są o 57--101\% gorsze
          od BKS. Instancje C mają wąskie okna czasowe -- algorytm ze słabą
          karą za spóźnienie nie jest wystarczająco motywowany do ich respektowania
          i generuje trasy o mniejszym dystansie kosztem wielu naruszeń.
    \item \textbf{Instancje R i RC:} nasze wyniki są \emph{poniżej} BKS (do $-23\%$).
          Jest to bezpośredni dowód naruszania okien czasowych -- minimalna
          odległość dopuszczalnych tras dla R101 wynosi 1650{,}80, więc wynik
          1265 jest fizycznie niemożliwy bez naruszeń. Mała kara (0{,}01) sprawia,
          że algorytm ignoruje okna czasowe i minimalizuje tylko dystans.
\end{itemize}

Podsumowanie: zastosowana penalizacja jest zbyt mała, aby wymusić dopuszczalność.
Porównanie bezpośrednie z BKS jest dlatego niemożliwe. Wartością eksperymentu jest
porównanie \textit{względne} ops\_on vs ops\_off w identycznych warunkach.

\subsection{Ograniczenia badania}

\begin{itemize}
    \item Zastosowana penalizacja (współczynnik 0{,}01 za spóźnienie) jest zbyt
          mała, by zapewnić dopuszczalność rozwiązań. Wyniki R/RC poniżej BKS
          jednoznacznie wskazują na naruszenia okien czasowych.
    \item Liczba powtórzeń $N=5$ daje ograniczoną moc statystyczną -- różnice
          poniżej 2\% mogą nie być istotne statystycznie.
    \item Strojenie metodą one-factor-at-a-time nie gwarantuje globalnego optimum
          w przestrzeni parametrów -- interakcje między parametrami nie są badane.
\end{itemize}

\end{document}
